\section{Introduction}
\label{sec:intro}

\emph{In one sentence:} under matched call budgets, six-vote self-consistency and six-paraphrase ensembling both deliver small, reliable out-of-sample gains over the single-sample expectation, the two sources are statistically indistinguishable from each other, and the apparent prompt-space margin is conditional --- largest on quantitative tasks, where the canonical instruction under-elicits and a single well-chosen instruction variant already approaches the ensemble.

Large language models answer questions through a single forward pass whose output depends on two design choices that the user controls: the prompt that specifies the task, and the decoding configuration that maps model probabilities into text. Both choices turn out to be fragile. A long line of work on prompt sensitivity shows that semantically equivalent reformulations of the same question can change model answers by large margins \citep{mizrahi2024state,sclar2024quantifying,leidinger2023language,zhu2023promptbench}. A parallel line of work on test-time compute shows that a single greedy decode is an unreliable summary of what the model knows, and that aggregating several decodes of the same prompt, as in self-consistency \citep{wang2023selfconsistency}, improves reasoning accuracy at the cost of additional inference calls.

These two observations motivate two ways of spending a fixed budget of model calls. The first spends the budget on the answer space: keep the prompt fixed, sample several decodes at nonzero temperature, and aggregate them by majority vote. The second spends the budget on the prompt side: present several semantically equivalent paraphrases of the task and aggregate the resulting answers; each paraphrase call is decoded stochastically at the same temperature as a resampled call, so this source mixes prompt variation with per-call decoding noise rather than isolating the prompt effect. Both strategies are used in practice, both appear in the evaluation literature under separate headings, and both consume the same scarce resource. Surprisingly, the literature offers little guidance on which source yields more accuracy per call, whether the two sources detect the same errors, and how the answer changes with model capability and task difficulty.

The question is becoming more urgent as model behavior shifts. Post-trained flagship models are increasingly confident: on the legacy benchmarks in our pool they answer near ceiling, and repeated sampling at moderate temperature often returns identical answers, which would render answer-space ensembling inert on exactly those benchmarks. At the same time, uncertainty estimation methods built on sampling, such as semantic entropy \citep{kuhn2023semantic,farquhar2024detecting}, implicitly assume that the sample distribution carries information about correctness. If that assumption fails for confident models, practitioners need to know before they build on it.

This paper presents, to our knowledge, the first budget-matched, within-question comparison of answer-space and paraphrase ensembling across multiple model families and task types; unlike prior prompt-ensembling work, which varies prompts but does not match call budgets against resampling on identical questions, our protocol holds the question, the model, and the total call budget constant by construction. Every question is answered sixteen times by every model under matched conditions: one greedy decode, six stochastic decodes of the original prompt, one stochastic decode of each of six validated paraphrases, and three additional second decodes that support a hybrid condition. From this single call matrix we derive self-consistency ensembles of any size up to six, paraphrase ensembles of the same sizes, a hybrid allocation, a combined twelve-vote ensemble, per-question agreement scores for both sources, and fragility measures that record how often paraphrasing alone flips the majority answer. The evaluation ran in two phases: a \emph{development suite} of 90 questions on six systems, whose unevenly truncated coverage (12 to 90 questions per model) reflects the free-tier quotas of the providers, fixed the conditions, strata, and rules; and a \emph{held-out suite} of 119 fresh questions, fully covered by four models from two providers and two families, which carries the primary evaluation. A greedy-paraphrase control that removes decoding stochasticity from the paraphrase side was collected on the development suite only, so out-of-sample paraphrase margins are read as properties of paraphrase ensembles \emph{with} sampling noise mixed in, not of pure prompt-space effects.

Our study addresses four research questions. First, under equal budgets, does spending calls on paraphrases buy more, equal, or less accuracy than spending them on resampling? Second, how brittle are model answers to paraphrasing per se, and where do paraphrase-induced flips concentrate? Third, can agreement across paraphrases serve as an error signal in the way that agreement across samples does? Fourth, how do these effects interact with model capability and task difficulty? The last question governs the first three: on the held-out suite, where the votes disagree, both sources deliver small reliable gains over the single-sample expectation (pooled $+2.4$ and $+4.9$ points, confidence intervals excluding zero) yet remain statistically inseparable from each other, and where the paraphrase-ensemble margin appears it concentrates on quantitative tasks whose canonical instruction under-elicits --- a conditional answer, not a universal winner. An instruction-variant analysis on the same data shows that this concentration coincides with an instruction effect rather than a pure aggregation effect: the best single instruction variant, selected post hoc, reaches $95.0\%$ single-call accuracy on SAT-Math against $91.2\%$ for the six-variant paraphrase ensemble, so the ensemble captures most, but not all, of what a well-chosen instruction achieves in a single call. Whether six resamples under the single best instruction would aggregate a further margin on top of it is a control we did not run and flag explicitly as such.

The empirical picture that emerges has a clear regime structure. When a model finds a task easy, its answer distribution collapses: on saturated tasks, models return the same answer for essentially every sample and every paraphrase, self-consistency and paraphrase ensembling both collapse onto the single-call baselines, and consistency scores carry no signal about which answers are wrong. When a task genuinely challenges a model, the votes disagree, and the interesting quantity stops being whether to aggregate and becomes what the aggregation is measured against: a single stochastic draw is already close to what six or twelve votes return on average, so the choice of the $k{=}1$ reference, greedy or single-sample, governs whether ensembling looks like a large gain or a wash, while the sign of the residual gain is set by whether the model's modal answer is correct.

Our contributions are as follows. We formalize test-time ensembling as a choice between two diversity sources under a shared budget, and we introduce a call matrix design that derives all conditions of interest from one matched set of calls, with development-suite controls that separate paraphrase variation from decoding stochasticity, answering mode, and task difficulty. We provide a multi-model, multi-benchmark empirical map of when each source helps, validated on a held-out suite of 119 fresh questions with four fully covered models, including evidence that paraphrase agreement functions as an uncertainty signal whose error coverage only partially overlaps that of sample agreement. We document the answer degeneracy phenomenon that renders both ensembling and consistency-based uncertainty estimation inert on saturated tasks for confident models. Finally, we translate the findings into an adaptive allocation rule checked by simulation on the development matrices and transferred to held-out questions, and we release the complete prompt materials, validated paraphrases, per-call predictions, and the scripts that recompute every table from the raw call logs in the repository accompanying this submission.

The remainder of the paper is organized as follows. Section~\ref{sec:related} reviews related work. Section~\ref{sec:method} formalizes the ensembling framework and the protocol. Section~\ref{sec:setup} describes models, benchmarks, and validation procedures. Section~\ref{sec:heldout} presents the held-out evaluation that carries the primary comparisons, with the development-suite analyses that motivated the design relegated to \ref{app:dev}; the greedy-paraphrase, hard-task, and chain-of-thought controls appear in \ref{app:greedypara}--\ref{app:cot}. Section~\ref{sec:discussion} discusses implications and limitations, and Section~\ref{sec:conclusion} concludes.
